\scanpage{47}\mspage{46}
% slanted
\margin{Ainsi, $\unsure{\mathcal M|M}$ s'identifie à un ens de parties de $\unsure{\mathcal F|\mathscr F}$ (savoir les $F_X$) i.e. à l'ens des figures élémentaires. \deleted{Donc on peut considérer que la structure est \unsure{portée} par $\mathcal F$, muni de $\mathcal F_{\unsure{\mathrm{él}|\mathrm{el}}}$ et de $\unsure{\Sigma}$, la relation de subdivision (relation d'ordre)} \deleted{Il y a \unsure{une sorte de symétrie|deux sortes de symétries} -- on travaille sur $\mathcal F$, \unsure{ou bien|comme} \inserted{par \unsure{les axiomes des}} \unsure{et tantôt} sur $\mathcal M$,  1°) relation d'ordre $\unsure{C'|\leq}$ (subdivision, ou $\ll$ raffinement) et ens. de parties (élémentaires, les figures).}}
\deleted{\unsure{Cela étant}}

\deleted{$\unsure{X}\ll\unsure{Y}\quad\unsure{F'}\unsure{\supset|\subset}F_X$} Si $F$ est une figure, alors $X\in F_{\deleted{\unsure{X|x}}}$ équivaut à $F_X\subset F$, qui implique que pour tout $Y\in F_X$ \inserted{i.e. $Y\leq X$} on a $F_Y\subset F$ i.e. $Y\in F$. Ainsi, $F$ est une partie de $\mathcal M$ \uline{fermée} pour $\leq$.

On dit que deux figures \inserted{$F,G$} \uline{sont compatibles} si $F\cup G$ est une figure, et que $X,Y\in\mathcal M$ sont \uline{compatibles} si $F_X,F_Y$ le sont

\uline{C \deleted{\unsure{2|3}}\unsure{3|2}} Pour que $F$ et $G$ soient compatibles, il f. et s. que $\forall\,\unsure{x|X}\in F,\ Y\in G$, $X$ et $Y$ le \unsure{soient|sont}.
% Two vertical emphasis strokes.
% slanted
\margin{correspond à Fig \unsure{4|9}}

\uline{Cor} \deleted{En particulier si $X$} Pour que $X,Y\in\mathcal M$ soient compatibles, il f et s. qu'ils \uuline{app}artiennent à une figure

\deleted{\uline{C 3} Si $F$ est.}
% slanted
\margin{\tikz[baseline=(a.base)]{\node[inner sep=1pt] (a) {\unsure{intervertir}}; \draw[<->] (a.north east) .. controls +(-.7,.2) and +(-.7,-.2) .. (a.south east);}}

\uline{C \deleted{\unsure{3|2}}\unsure{2|3}} Soit $F$ une figure, alors tte partie de $F$ fermée pour $\leq$ est aussi une figure (i.e. pour qu'une partie $F'$ de $F$ soit une figure, il f et il \uline{suffit} qu'elle soit fermée pour $\leq$).
% Two vertical emphasis strokes.

(\uline{Sous-figures}) Si $\mathcal M\ne\varnothing$ cela \deleted{\unsure{signifie}} \unsure{implique} que la partie vide de $\mathcal M$ est une figure. Si $\mathcal M=\varnothing$, \deleted{cela implique que} \inserted{on aura $\mathcal F=\varnothing$} \deleted{si $\mathcal F$} ou $\mathcal F=\{\varnothing\}$, et on suppose que c'est le deuxième cas : \uline{figure vide} existe tjrs
% slanted
\margin{\uline{Contrée vide} si $\mathcal M=\varnothing$}
